\section*{Erweiterte Eigenständigkeitserklärung zur Verwendung generativer KI-Systeme als Hilfsmittel}\label{sec:extended_doi}


\vspace{.25cm}

\begin{framed}
\noindent\textbf{Hiermit versichere ich, dass ich diese Prüfungsleistung selbständig verfasst habe. Die verwendeten Quellen und Hilfsmittel habe ich angegeben.}
\end{framed}

Sofern ich generative KI-Systeme über die Prüfung von Grammatik und Rechtschreibung von eigenständig erstelltem Text hinaus genutzt habe, bestätige ich hiermit, dass ich: 
\begin{itemize}
  \item die von generativen KI-Systemen direkt übernommenen Inhalte im inhaltlichen Teil der Arbeit entsprechend markiert bzw. zitiert habe, 
  \item sichergestellt habe, dass die mit Hilfe der genannten generativen KI-Systemen erzeugten und sonstigen von mir verwendeten Inhalte faktisch korrekt sind,
  \item mir darüber im Klaren bin, dass ich als Verfasserin oder Verfasser dieser Arbeit die Verantwortung für die darin enthaltenen Angaben und Aussagen übernehme,  
  \item sofern ich generative KI-Systeme zur Unterstützung wesentlicher Arbeitsschritte eingesetzt habe, die nicht bereits im inhaltlichen Teil der Arbeit oder den Referenzen thematisiert bzw. zitiert werden, dies tabellarisch erfasst habe. Die Tabelle habe ich mit Abgabe der Arbeit eingereicht (siehe Anhang). 
\end{itemize}

\vspace{3cm}
\noindent Unterensingen, 29.08.2025,\\

\noindent\textbf{Ort, Datum, Unterschrift}

\newpage

\subsection*{Anhang der erweiterten Eigenständigkeitserklärung}

\renewcommand{\arraystretch}{1.5} 
\newcolumntype{L}[1]{>{\raggedright\arraybackslash}p{#1}}
\newcolumntype{C}[1]{>{\centering\arraybackslash}p{#1}}

\begin{table}[htbp]
  \centering
  \begin{tabularx}{\textwidth}{
    @{} 
      L{3.5cm}
      L{2.5cm}
      C{2.5cm}
      L{5cm}
    @{}}
    \toprule
    \textbf{Arbeitsschritt}
      & \textbf{Eingesetzte Gen‑KI‑ Systeme}
      & \textbf{Verwendung in Kapitel}
      & \textbf{Beschreibung der Verwendungsweise} \\
    \midrule
    Sprachliche Optimierung und Strukturierung von Text
      & DeepL Write, Grammarly, Google Gemini (\textit{2.5‑pro})
      & Alle
      & Verbesserung der Sprache, Rechtschreibung und Strukturierung der einzelnen Kapitel. \\
    Zusammenfassung von Literatur
      & OpenAI ChatGPT (\textit{o3‑mini‑high})
      & 3
      & Zusammenfassung wissenschaftlicher Texte zur initialen Bewertung ihrer Relevanz. \\
    Generierung von Docstrings und Kommentaren
      & Google Gemini (\textit{2.5‑pro})
      & Code
      & Erstellung von Kommentaren und Docstrings im Source‑Code zur Steigerung der Lesbarkeit. \\
    \bottomrule
  \end{tabularx}
  \caption{Einsatz generativer KI‐Systeme in der Arbeit}
\end{table}